\section*{Chapitre \ref{ch:b2:thermal-radiation} --- Rayonnement thermique}

\begin{solution}{exo:b2:thermal-radiation:1}
$\lambda_{\max}$~: \qty{0.50}{\micro m}, \qty{9.4}{\micro m}, \qty{1.04}{\micro m},
\qty{1.06}{mm}, \qty{38}{\micro m}, \qty{3.2}{\micro m}~; $\sigma T^4$~: \qty{6.4e7}{},
\qty{520}{}, \qty{3.5e6}{}, \qty{3.1e-6}{}, \qty{2.0}{}, \qty{3.7e4}{W/m^2}. Visibles~:
le Soleil et le filament~; le tisonnier rougeoie sourdement grâce à la
queue de courtes longueurs d'onde, un dix-millième de sa puissance.
\end{solution}

\begin{solution}{exo:b2:thermal-radiation:2}
$S = \sigma T_\odot^4(R_\odot/d)^2 = \qty{1360}{W/m^2}$~; $L = 4\pi d^2S =
\qty{3.8e26}{W}$~; $\pi R^2S = \qty{1.7e17}{W}$~; \qty{2e7}{W} par personne --- neuf
mille fois la consommation mondiale.
\end{solution}

\begin{solution}{exo:b2:thermal-radiation:3}
Émis \qty{830}{W}, absorbés \qty{700}{W}, net \qty{130}{W}~; $h_{\text{r}} =
4\varepsilon\sigma T_0^3 = \qty{5.6}{W/m^2/K}$. Plus que le métabolisme~:
les vêtements le réduisent~; des murs à \qty{15}{\celsius} ajoutent quelque \qty{50}{W} de perte
--- on a froid dans de l'air chaud entre des murs froids.
\end{solution}

\begin{solution}{exo:b2:thermal-radiation:4}
$A = P/\varepsilon\sigma T^4 = \qty{49}{mm^2}$~; $\ell = A/\pi d = \qty{31}{cm}$~;
\qty{4.8}{W} de lumière~; une LED de \qty{16}{W}.
\end{solution}

\begin{solution}{exo:b2:thermal-radiation:5}
(a) $\eu^x - 1 \approx x$~: $u_\nu = (8\pi\nu^2/c^3)\,k_BT$ --- $k_BT$ par mode.
(b) $\int\nu^2\dd\nu$ diverge. (c) $\nu = k_BT/h = \qty{6e12}{Hz}$ (\qty{50}{\micro m})~:
les radiofréquences et les micro-ondes à \qty{300}{K} sont bien dans le régime classique (le
bruit d'une résistance vaut $k_BT$ par unité de bande). (d) La probabilité de
trouver le quantum $h\nu$ excité à $T$~: un facteur de Boltzmann.
\end{solution}

\begin{solution}{exo:b2:thermal-radiation:6}
(a) $x = 4.965$. (b) $hc/4.965k_B = \qty{2.90e-3}{m.K}$. (c) $x^3/(\eu^x - 1)$
est maximal là où $3(1 - \eu^{-x}) = x$~: $2.82$~; les densités par unité de fréquence
et par unité de longueur d'onde sont des fonctions différentes ($\dd\nu = c\,\dd\lambda/\lambda^2$)
et culminent en des points différents. (d) \qty{10000}{K}~; blanc bleuté.
\end{solution}

\begin{solution}{exo:b2:thermal-radiation:7}
(a) $u = (8\pi^5k_B^4/15h^3c^3)T^4$, $\sigma = cu/4T^4 = \qty{5.67e-8}{W/m^2/K^4}$.
(b) L'isotropie et la projection en $\cos\theta$~: $\int\cos\theta\,\dd\Omega/4\pi = 1/4$
sur un hémisphère. (c) $n = 8\pi \times 2.404\,(k_BT/hc)^3 = 60.4\,(k_BT/hc)^3$~;
énergie moyenne $(\pi^4/15)/2.404 = 2.70\,k_BT$. (d) $5 \times 10^8$ et $400$
par centimètre cube.
\end{solution}

\begin{solution}{exo:b2:thermal-radiation:8}
(a) $\sigma(T_1^4 - T_2^4)$. (b) $T_{\text{s}}^4 = (T_1^4 + T_2^4)/2$~; moitié moins. (c)
$1/(n + 1)$. (d) En sommant les réflexions, $\sigma(T_1^4 - T_2^4)/(2/\varepsilon -
1)$~; $\varepsilon = 0.05$~: $\qty{460}{W/m^2}/39 = \qty{12}{W/m^2}$ --- un vase de
\qty{0.05}{m^2} perd \qty{0.6}{W}, évaporant un quart de kilogramme
d'azote par jour.
\end{solution}

\begin{solution}{exo:b2:thermal-radiation:9}
(a) $\varepsilon\sigma(T^4 - 250^4) = h(278 - T)$~: $T \approx \qty{266}{K}$. (b)
\qty{277}{K}. (c) Un $h$ plus grand fixe la surface à l'air. (d) Le feuillage
à \qty{278}{K} remplace le ciel à \qty{250}{K}.
\end{solution}

\begin{solution}{exo:b2:thermal-radiation:10}
(a) \qty{231}{K}, \qty{255}{K}, \qty{210}{K}. (b) \qty{506}{K}, \qty{33}{K}, \qty{5}{K}.
(c) Chaque couche opaque rayonne vers le haut et vers le bas~; les bilans donnent
$T_k^4 = (k + 1)T_{\text{e}}^4$ depuis le sommet~: $n \approx 100$ pour Vénus. (d)
$((1 - A)S/\sigma)^{1/4} = \qty{381}{K}$~; sans atmosphère rien ne porte la
chaleur vers la face nocturne, et le régolithe emmagasine peu~: la surface
rayonne la chaleur de sa journée en un jour environ et se stabilise vers
\qty{100}{K}.
\end{solution}

\begin{solution}{exo:b2:thermal-radiation:11}
(a) $4\sigma T_{\text{e}}^3 = \qty{3.8}{W/m^2/K}$. (b) \qty{1.0}{K}. (c) $C =
\qty{2.1e8}{J/m^2/K}$, $\tau = C/4\sigma T_{\text{e}}^3 \approx \qty{5.6e7}{s}$, deux
ans. (d) Un air plus chaud contient plus de vapeur d'eau (un gaz à \omterm{prop:b2:thermal-radiation:greenhouse}{effet de serre}) et
fait fondre la glace (albédo plus bas)~: les deux amplifient le réchauffement initial.
\end{solution}

\begin{solution}{exo:b2:thermal-radiation:12}
(a)
\[
\frac{M_1}{M_2} = \Big(\frac{\lambda_2}{\lambda_1}\Big)^5\exp\Big[-\frac{hc}{k_BT}\Big(\frac{1}{\lambda_1} - \frac{1}{\lambda_2}\Big)\Big] :
\]
$\varepsilon$ se simplifie. (b) $T = \qty{14388}{\micro m.K} \times 0.427/(5\ln 1.385 +
\ln 10) = \qty{1570}{K}$. (c) Reçu $\propto 0.1 \times 350 + 0.9 \times 293 =
299$~; annoncé $(299 - 0.05 \times 293)/0.95 \approx \qty{299}{K}$~: elle manque les
\qty{350}{K}. (d) Son \omterm{prop:b2:thermal-radiation:kirchhoff}{émissivité} est celle qu'on a supposée.
\end{solution}

\begin{solution}{pb:b2:thermal-radiation:1}
\textbf{1.} \qty{5800}{K}.

\textbf{2.} \qty{6.4e7}{W/m^2}~; \qty{3.9e26}{W}.

\textbf{3.} $L/4\pi d^2 = \qty{1.37e3}{W/m^2}$~; $u = S/c = \qty{4.6e-6}{J/m^3}$~;
$P = S/c = \qty{4.6e-6}{Pa}$ si absorbée, le double si réfléchie.

\textbf{4.} $2.7k_BT = \qty{2.2e-19}{J}$ (\qty{1.35}{eV})~; $6 \times 10^{21}$
photons par seconde et par mètre carré.

\textbf{5.} $3.9 \times 10^{26} \times 1.4 \times 10^{17} = \qty{5.5e43}{J}$~; $Mc^2 =
\qty{1.8e47}{J}$~: trois dix-millièmes de sa masse.

\textbf{6.} \qty{1.2e17}{W}~; $T_{\text{e}} = \qty{255}{K}$.

\textbf{7.} Intercepté sur $\pi R^2$, émis depuis $4\pi R^2$ grâce à la
rotation et aux vents~; le seul point subsolaire~: $((1 - A)S/\sigma)^{1/4} =
\qty{360}{K}$.

\textbf{8.} \qty{381}{K}~; la nuit la chaleur emmagasinée ($\sim C\Delta T$) s'écoule par
rayonnement avec la constante de temps $C/4\sigma T^3 \approx \qty{1}{jour}$, courte
devant la quinzaine d'obscurité~: le sol tombe là où son
rayonnement égale le filet venu d'en dessous, \qty{100}{K}.

\textbf{9.} \qty{10}{\micro m}, vingt fois celle du Soleil.

\textbf{10.} $(0.65/0.70)^{1/4} = 0.982$~: \qty{4.7}{K} de moins, \qty{250}{K}.

\textbf{11.} Espace~: $\sigma T_{\text{a}}^4 = \sigma T_{\text{e}}^4$~; couche~: $\sigma
T_{\text{s}}^4 = 2\sigma T_{\text{a}}^4$~; surface~: $\sigma T_{\text{s}}^4 = \sigma
T_{\text{e}}^4 + \sigma T_{\text{a}}^4$. D'où $T_{\text{s}} = 2^{1/4}T_{\text{e}} = \qty{303}{K}$.

\textbf{12.} $T_{\text{s}}^4 = T_{\text{e}}^4/(1 - \varepsilon/2)$~; $(288/255)^4 = 1.63$
donne $\varepsilon = 0.77$.

\textbf{13.} La vapeur d'eau, le dioxyde de carbone, le méthane et l'ozone absorbent dans
l'infrarouge (bandes de vibration~: H$_2$O vers \qty{6}{\micro m} et au-delà de
\qty{17}{\micro m}, CO$_2$ à \qty{15}{\micro m})~; N$_2$ et O$_2$ non (pas de
dipôle)~; la fenêtre est \qty{8}{}--\qty{13}{\micro m}.

\textbf{14.} $T_{\text{a}}^4 = T_{\text{s}}^4/2$~: $\varepsilon\sigma T_{\text{a}}^4 \approx
\qty{150}{W/m^2}$, contre \qty{238}{W/m^2} de lumière solaire absorbée (une
atmosphère à plusieurs couches en renvoie davantage, quelque \qty{330}{W/m^2}).

\textbf{15.} $3.7/3.8 \approx \qty{1}{K}$.

\textbf{16.} \qty{2.1e8}{J/m^2/K}~; $\tau \approx \qty{5.6e7}{s}$, deux ans.

\textbf{17.} Ils réfléchissent la lumière solaire et rayonnent de l'infrarouge vers le bas~; la nuit
seul le second agit~: les nuits nuageuses sont plus douces.

\textbf{18.} \qty{49}{mm^2}~; \qty{31}{cm}~; enroulé pour tenir dans l'ampoule et pour réduire
les pertes par convection vers le gaz de remplissage.

\textbf{19.} \qty{1.04}{\micro m}~; \qty{4.8}{W}~; $8\%$.

\textbf{20.} $R = V^2/P = \qty{880}{\ohm}$~; à froid \qty{88}{\ohm}~: \qty{2.6}{A}
au lieu de \qty{0.26}{A} --- les ampoules lâchent à l'allumage.

\textbf{21.} $49 \times (2800/3200)^4 = \qty{29}{mm^2}$~; \qty{8.4}{W} de visible~; huit
doublements de l'évaporation, une durée de vie $250$ fois plus courte --- des heures.

\textbf{22.} Le tungstène évaporé se lie à l'halogène et est rendu
au filament chaud au lieu de noircir le verre~: un filament plus chaud
à durée de vie normale, dans une petite enveloppe de quartz.

\textbf{23.} \qty{55}{W} d'infrarouge et de conduction par le gaz~; le
verre absorbe une part de l'infrarouge.

\textbf{24.} $C = \qty{2.8e-3}{J/K}$~; avec $P \propto T^4$, le temps pour que $T$ soit divisée par deux
vaut $7CT_0/3P \approx \qty{0.3}{s}$ --- et l'incandescence suit le secteur
avec un retard de \qty{0.04}{s}, d'où le fait qu'elle ne papillote guère.

\textbf{25.} La température fixe la forme et le maximum du spectre
($\lambda_{\max}T$), et le total ($\sigma T^4$)~; le filament, le Soleil et la Terre
sont trois \omterm{def:b2:thermal-radiation:blackbody}{corps noirs}, à \qty{2800}{}, \qty{5800}{} et \qty{288}{K}, chacun
équilibrant ce qu'il reçoit par $\sigma T^4$.
\end{solution}
